\documentclass[11pt,a4paper]{article}

\usepackage[margin=24mm]{geometry}
\usepackage[T1]{fontenc}
\usepackage[utf8]{inputenc}
\usepackage{lmodern}
\usepackage{microtype}
\usepackage{amsmath}
\usepackage{xcolor}
\usepackage{booktabs}
\usepackage{tabularx}
\usepackage{array}
\usepackage{enumitem}
\usepackage{xurl}
\usepackage[hidelinks]{hyperref}
\emergencystretch=2em
\newcolumntype{Y}{>{\raggedright\arraybackslash}X}
\newcolumntype{L}[1]{>{\raggedright\arraybackslash}p{#1}}

\definecolor{pilotblue}{HTML}{2368FF}
\definecolor{pilotgreen}{HTML}{16834A}
\definecolor{pilotamber}{HTML}{B56A00}
\definecolor{pilotred}{HTML}{C52F46}
\newcommand{\Pilot}{\textbf{Pilot}}
\newcommand{\Fabric}{\textbf{AION Device Fabric}}
\newcommand{\implemented}{\textcolor{pilotgreen}{\textbf{Implemented}}}
\newcommand{\bounded}{\textcolor{pilotamber}{\textbf{Bounded}}}
\newcommand{\deferred}{\textcolor{pilotred}{\textbf{Production work}}}

\title{\textbf{Pilot Live Companion, Screen Understanding,\\Provider Evidence, and Observer Mode}\\[4pt]
\large AION Device Fabric Technical Supplement}
\author{Tessaris / AION Product and Technical Record}
\date{29 August 2026\\Fabric implementation versions 0.15.0--0.19.0}

\begin{document}
\maketitle

\begin{abstract}
This supplement documents the implemented transition from television control
to evidence-grounded television understanding in AION Device Fabric.  Fabric
0.15 introduced Pilot Live Companion and signed, rights-aware Pilot Moments.
Fabric 0.16 added local phone-camera perception, positioned optical character
recognition (OCR), structured screen fusion, confidence-labelled findings, and
private owner correction.  Fabric 0.17 added provider-reference normalization,
optional official YouTube metadata enrichment, and a short-lived explicit
Observer Mode authority.  Together these capabilities allow Pilot to reason
about recent television speech and owner-supplied visual evidence without
claiming unsupported access to protected programme pixels, copying protected
media, inferring a person's identity from a face, or treating provider metadata
as proof of subscription entitlement.  This document records the architecture,
schemas, execution flows, trust boundaries, verification results, and remaining
production work.
\end{abstract}

\section{Scope and Implemented Status}

The work described here is additive to the existing mother-brain, signed
capability, webOS gateway, private-phone, audit-ledger, and television-Canvas
architecture.  It does not replace COMDEX intelligence files and does not grant
the television direct access to COMDEX memory, service credentials, device
pairing secrets, or private household identity data.

\begin{table}[htbp]
\centering
\small
\caption{Implemented releases covered by this supplement.}
\label{tab:pilot-live-releases}
\begin{tabularx}{\textwidth}{@{}L{14mm}L{34mm}Y L{25mm}@{}}
\toprule
Release & Principal object & Capability introduced & Verification\\
\midrule
0.15.0 & \nolinkurl{pilot.live.context.v1} and
\nolinkurl{pilot.moment.v1} & Explicit live fact-check, explanation, and
rights-aware Moment preparation using recent memory-only audio and permitted
TV state & 136 tests; real LG state observation and Moment preparation\\
0.16.0 & \nolinkurl{pilot.screen-understanding.v1} & Local Apple Vision OCR
geometry, screen-element detection, source/confidence fusion, and persona-bound
correction & 144 tests; real OCR geometry processing; no live TV commands\\
0.17.0 & \nolinkurl{pilot.provider-metadata.v1} and
\nolinkurl{pilot.explicit-observer.v1} & Provider identifiers, optional official
YouTube enrichment, and explicit short-lived camera-session leases & 149 tests;
audit chain valid; no live TV commands\\
0.18.0 & \nolinkurl{pilot.live-news.v1} & Claim decomposition, source-labelled
comparison, correction timeline, playback-preserving spoken verdict, and
private-phone evidence record & 154 tests; safe fallback verified; no live TV
commands\\
0.18.1 & Evidence-provider failover & OpenAI Responses web search followed by
Gemini Google Search grounding and then an explicit non-verdict handoff & 156
tests; live grounded failover verified; no live TV commands\\
\bottomrule
\end{tabularx}
\end{table}

The user-facing commands introduced by this programme include:

\begin{itemize}[leftmargin=*]
  \item ``Pilot, fact check that.''
  \item ``Pilot, explain that.''
  \item ``Pilot, share that moment.''
  \item ``Pilot, what is on the screen?''
  \item ``Pilot, what is the score?''
  \item ``Pilot, read the subtitle.''
  \item ``Pilot, what product is that?''
\end{itemize}

These phrases are entry points into governed intents rather than authority to
perform arbitrary external actions.  A request may produce an answer, a
prepared context object, an explicit request for more evidence, or an honest
refusal.  It does not silently purchase, message, publish, identify a person,
or export protected media.

\section{System Architecture}

\subsection{Evidence flow}

The implemented data path is:

\begin{center}
\begin{tabular}{c}
Television programme, game, browser, or application\\
$\downarrow$ permitted metadata \hspace{12mm} $\downarrow$ room audio
\hspace{12mm} $\downarrow$ explicit phone capture\\
webOS gateway \hspace{9mm} memory-only audio ring
\hspace{9mm} local Apple Vision\\
$\searrow$ \hspace{27mm} $\downarrow$ \hspace{27mm} $\swarrow$\\
AION mother-brain evidence normalization and policy\\
$\downarrow$\\
LiveContext $+$ provider metadata $+$ OCR geometry $+$ labels $+$ screen belief\\
$\downarrow$\\
confidence-labelled screen-understanding object\\
$\downarrow$\\
TV Canvas answer, private phone detail, correction, or governed next action
\end{tabular}
\end{center}

The restricted television remains a gateway-managed surface.  The mother node
performs planning, canonicalization, signing, local perception, policy checks,
evidence retrieval, and auditing.  The private phone is both a control surface
and an optional owner-visible sensor.  The TV Canvas receives only a sanitized
projection.

\subsection{Evidence classes}

Pilot distinguishes rather than silently merges its inputs:

\begin{enumerate}[leftmargin=*]
  \item \textbf{Device evidence:} permitted webOS foreground-application,
        audio, playback, and cached state.
  \item \textbf{Ephemeral transcript evidence:} locally transcribed material
        from the bounded recent-audio buffer, used only after an explicit
        request.
  \item \textbf{Owner-captured visual evidence:} OCR observations, bounding
        boxes, classifications, and an image digest derived from an explicit
        phone submission.
  \item \textbf{Provider evidence:} normalized official URLs or identifiers
        and, where configured, bounded official API metadata.
  \item \textbf{Prior screen belief:} the most recent confidence-labelled
        Fabric view of the active surface.
  \item \textbf{Owner correction:} a persona-bound statement correcting one
        supported field of the current observation.
\end{enumerate}

Every fusion result records source descriptions and confidence values.  A
source conflict is represented explicitly.  Missing evidence is not converted
into fabricated certainty.

\section{Fabric 0.15: Pilot Live Companion}

\subsection{Memory-only recent-audio window}

The voice service maintains a bounded in-memory ring containing at most
approximately 30 seconds of recent audio.  In the current implementation the
ring is a fixed-length deque sampled in bounded blocks.  It is cleared when the
microphone sleeps and when the voice process stops.  Raw audio is not written
to the Fabric runtime directory, capability ledger, Canvas, phone interface,
or Moment object.

The ring is not itself an authorization to retain or research television
speech.  Persistence begins only when the user makes an explicit contextual
request.  The mother then constructs a small structured LiveContext from
bounded recent transcript fragments and freshly observed permitted TV state.

\subsection{LiveContext schema}

The persisted object has schema identifier
\texttt{pilot.live.context.v1}.  Its principal fields are:

\begin{table}[htbp]
\centering
\small
\begin{tabularx}{\textwidth}{@{}L{37mm}Y@{}}
\toprule
Field & Meaning\\
\midrule
\texttt{created\_at} & UTC creation time.\\
\texttt{surface} & Current confidence-labelled Fabric surface.\\
\texttt{app} & Bounded foreground application identifier and title.\\
\texttt{playback} & Exposed title, state, position, or duration fields when
available.\\
\texttt{recent\_statement} & Most recent bounded, non-wake transcript fragment.\\
\texttt{recent\_transcripts} & At most four normalized transcript fragments.\\
\texttt{evidence} & Typed descriptions of webOS and ephemeral transcript
sources.\\
\texttt{confidence} & Context confidence, not a truth score for the statement.\\
\texttt{privacy} & Raw-audio retention and explicit-request declarations.\\
\texttt{context\_hash} & Canonical hash of the structured record.\\
\bottomrule
\end{tabularx}
\end{table}

The context hash is calculated over canonical bytes.  It permits later audit
records and Moment objects to refer to the exact evidence object without
embedding raw audio.

\subsection{Fact-check and explanation flow}

For ``fact check that'' or ``explain that'', the runtime performs the following
bounded sequence:

\begin{enumerate}[leftmargin=*]
  \item Verify the signed capability corresponding to the requested intent.
  \item Observe fresh permitted webOS state and record its action receipt.
  \item Extract a bounded recent statement from the memory-only transcript
        window; if unavailable, use a sufficiently supported captured subtitle
        or owner correction.
  \item Construct and hash the LiveContext.
  \item Route the question to the research layer with the appropriate
        fact-check or explanation mode.
  \item Prefer current, attributable evidence and retain sanitized results
        rather than the raw provider response.
  \item Present a compact result on the TV Canvas and fuller evidence on the
        private phone.
  \item Write an audit event containing the intent, evidence hashes, and result
        state.
\end{enumerate}

A fact-check result is evidence-labelled.  Where research evidence is not
available, Pilot must not invent a verdict; it returns an insufficient-evidence
state or an explicit search handoff.  The system therefore distinguishes
\emph{what was said}, \emph{what the TV reports is active}, and \emph{what
external evidence supports}.

\section{Fabric 0.15: Signed Pilot Moments}

\subsection{Purpose and rights boundary}

``Share that moment'' currently prepares a portable context object rather than
copying television media.  The object uses schema
\texttt{pilot.moment.v1}.  It may contain application and playback metadata, a
bounded recent statement, the originating context hash, a requested interval
of at most 30 seconds, and an optional recipient hint.  It explicitly records:

\begin{itemize}[leftmargin=*]
  \item \texttt{protected\_audio\_copied = false};
  \item \texttt{protected\_video\_copied = false};
  \item provider-native sharing is preferred; and
  \item a fallback card contains context rather than programme media.
\end{itemize}

The current provider-aware hierarchy is:

\begin{enumerate}[leftmargin=*]
  \item YouTube-like surfaces may prepare a provider timestamp link.
  \item Twitch-like surfaces may prepare a provider-native clip operation,
        pending an authorized provider adapter.
  \item Other surfaces produce a signed context card unless an official native
        sharing mechanism is exposed.
\end{enumerate}

The Moment remains \texttt{prepared\_not\_sent}.  Selecting a recipient and
sending through a messaging service remain separate private actions.

\subsection{Integrity model}

Let $M$ be the canonical Moment payload.  Fabric calculates

\begin{equation}
 h_M = \operatorname{SHA256}(\operatorname{canonical}(M)),
\end{equation}

and signs the canonical bytes with the issuing node identity:

\begin{equation}
 \sigma_M = \operatorname{Sign}_{K_{\mathrm{issuer}}}
             (\operatorname{canonical}(M)).
\end{equation}

The stored record includes $M$, $h_M$, the issuer public key, and $\sigma_M$.
This establishes provenance and tamper evidence; it does not create media
rights or prove that a recipient is entitled to play the referenced programme.

\section{Fabric 0.16: Non-Disruptive Screen Understanding}

\subsection{Owner-visible phone perception}

Because LG webOS does not expose a supported general-purpose third-party
screen-capture interface, the implemented visual path uses a photograph
submitted explicitly from the private phone.\cite{webosscreenshot} The service
accepts JPEG, PNG, or HEIC content up to 5 MiB.  It writes a
permission-restricted temporary file,
runs local Apple Vision, and deletes the temporary image in a
\texttt{finally} cleanup path.

The persisted \texttt{aion.phone.visual-perception.v1} record contains only:

\begin{itemize}[leftmargin=*]
  \item SHA-256 digest and byte count of the submitted image;
  \item explicit-capture source and consent labels;
  \item at most 40 normalized OCR strings;
  \item at most 40 positioned OCR observations;
  \item at most 12 bounded image-classification labels;
  \item inferred surface, view, cues, confidence, and timestamp; and
  \item \texttt{image\_retained = false}.
\end{itemize}

Each positioned OCR item contains bounded normalized coordinates
$x,y,w,h\in[0,1]$ and recognition confidence.  Geometry allows Pilot to
distinguish likely subtitle regions from menus and other screen text.

\subsection{Fusion model}

The screen fusion service combines phone perception $P$, permitted device state
$D$, explicit live context $L$, previous screen belief $B$, and provider
metadata $R$:

\begin{equation}
 S = F(P,D,L,B,R).
\end{equation}

The current implementation is a bounded deterministic evidence classifier,
not an unrestricted claim of visual understanding.  For each supported finding
$j$, it emits detection state $d_j$, confidence $c_j$, supporting source IDs,
and a normalized evidence subset.  The overall confidence is the mean of
detected supported-element confidences, or the base perception confidence when
no supported element is detected:

\begin{equation}
 C(S)=
 \begin{cases}
  |J|^{-1}\sum_{j\in J}c_j, & |J|>0,\\
  c_P, & |J|=0.
 \end{cases}
\end{equation}

This value indicates evidence strength within the implemented detector set; it
is not a probability that every semantic statement about the scene is true.

\subsection{Implemented detectors}

\begin{table}[htbp]
\centering
\small
\caption{Current screen-understanding findings.}
\begin{tabularx}{\textwidth}{@{}L{28mm}Y Y@{}}
\toprule
Finding & Evidence used & Important boundary\\
\midrule
Subtitles & OCR position, centre alignment, line length, and recognition
confidence & Likely subtitle text, not guaranteed programme transcript\\
Scoreboard & Team/score pattern, match clock, sports terms, and image labels &
Does not independently verify the official match result\\
Products & Currency/price expressions, purchase language, and product labels &
Any save, basket, or purchase remains private-confirmation-only\\
Locations & Map, route, distance, station, airport, landmark, and city cues &
Does not infer a private location or authorize navigation\\
Game HUD & Game surface state plus mission, level, score, health, or quest text &
Does not control gameplay or infer hidden game state\\
People & Generic person-related image labels & Faces remain unknown; identity is
never guessed\\
\bottomrule
\end{tabularx}
\end{table}

The resulting \texttt{pilot.screen-understanding.v1} object contains an
observation ID, app and surface, summary, confidence, source-conflict flag,
typed source records, supported findings, provider metadata, corrections,
privacy declarations, and canonical evidence hash.

\subsection{Private correction}

The phone controller lets the owner correct only the latest observation and
only an allowlisted field: subtitle, scoreboard, product, location, or scene.
The correction is normalized, bounded to 240 characters, attached to a private
persona, timestamped, and incorporated into a new evidence hash.  A stale
observation ID is rejected.  Person identity is intentionally absent from the
allowlist.

Corrections are evidence supplied by the owner, not automatic ground truth.
They can support later explanations or fact-check requests, but should remain
source-labelled in any downstream result.

\section{Fabric 0.17: Provider Evidence}

\subsection{Normalized provider record}

The \texttt{ProviderMetadataResolver} accepts a bounded set of HTTPS references
obtained from the current evidence context.  It recognizes supported YouTube
video IDs, Netflix title IDs, and Twitch clip IDs and writes a
\texttt{pilot.provider-metadata.v1} object.  Common fields include provider,
content ID, canonical official URL, metadata status, confidence, source links,
limitations, application evidence, entitlement flags, and evidence hash.

The following invariants always apply:

\begin{itemize}[leftmargin=*]
  \item \texttt{availability\_verified = false} unless a future authorized
        availability adapter proves otherwise;
  \item \texttt{account\_entitlement\_verified = false} unless a future
        persona-specific account adapter proves otherwise;
  \item metadata is not treated as permission to play, copy, clip, or send
        protected content; and
  \item an unavailable provider API cannot prevent local screen perception.
\end{itemize}

\subsection{YouTube enrichment}

Without a configured API key, a recognized YouTube URL produces a
\nolinkurl{provider_identifier_only} record.  When the owner configures
\nolinkurl{YOUTUBE_API_KEY}, the mother may call the official YouTube Data API
\nolinkurl{videos.list} endpoint using a restricted field projection.  The
response is limited to 500,000 bytes and may enrich the record with title,
channel, publication time, ISO-8601 duration, caption availability, live state,
and embeddability.\cite{youtubevideos} A timeout, invalid key, oversized
response, invalid JSON,
or other provider failure produces \nolinkurl{official_api_unavailable} while
preserving the identifier-only evidence path.

\subsection{Netflix and Twitch boundaries}

A Netflix title reference records the official title ID and identifies
Netflix's supported mobile Moments feature as a potential native-share
route.\cite{netflixmoments}
It does not claim that Moments can be driven from the LG browser, that a title
is available in the household's region, or that Pilot may copy the programme.

A Twitch clip reference records the official clip identifier and a native-clip
boundary.  Enriched Twitch operations remain unavailable until a developer
application and user-authorized adapter are connected.  Twitch embedded video
also imposes HTTPS and parent-domain requirements that a production surface
must satisfy.\cite{twitchembed}

\section{Fabric 0.17: Explicit Observer Mode}

\subsection{Authority and session lifecycle}

Observer Mode is a backend and phone-interface foundation for repeated visual
sampling.  It is deliberately not a hidden camera mode.  The owner must press
\emph{Start Observer Mode} on the private phone and grant visible browser or
native camera permission.

Starting a session creates:

\begin{itemize}[leftmargin=*]
  \item a random session identifier;
  \item a 256-bit URL-safe bearer token returned once to the phone;
  \item the initiating private persona ID;
  \item a start and expiry time;
  \item a maximum duration of 600 seconds;
  \item a configured interval clamped to 3--30 seconds;
  \item accepted-frame counters; and
  \item declarations that source frames are not retained and visible owner
        start is required.
\end{itemize}

Only the token hash is persisted.  The token itself is absent from status
snapshots and audit events.  For an incoming frame with token $t$, the service
requires

\begin{equation}
 \operatorname{CTEqual}(\operatorname{SHA256}(t),h_t)=\mathrm{true},
\end{equation}

where \texttt{CTEqual} denotes constant-time comparison and $h_t$ is the
stored token hash.  It additionally requires an active unexpired session and a
time since the previous accepted frame greater than or equal to the permitted
interval.  Invalid, early, stopped, or expired frames are rejected before image
analysis.

Stopping the observer is bound to the same private persona.  It marks the
session inactive and clears the persisted token hash.  The start and stop
events enter the Fabric audit chain without the bearer token or image content.

\subsection{Phone implementation}

The controller requests an environment-facing camera stream with audio
disabled, shows the live preview, downsizes frames to at most 1,280 pixels on
the largest requested dimension, encodes bounded JPEG samples, and sends them
to the existing perception endpoint with the
\texttt{X-Pilot-Observer} session header.  The page stops its media tracks when
the user stops Observer Mode or leaves the page.

Each accepted frame follows the same local Apple Vision, immediate deletion,
screen-fusion, hash, and audit path as a manual capture.  Observer Mode changes
capture cadence; it does not bypass the normal perception or capability
boundary.

\subsection{Secure-context limitation}

Browser \texttt{getUserMedia()} is restricted to secure contexts.\cite{mdnmedia}
The current phone controller is served over token-protected local HTTP for the
developer proof.  On mobile browsers enforcing the secure-context rule, the UI
therefore refuses automatic observation and directs the owner to the existing
manual \emph{Take a picture of the TV screen} control.  The implementation does
not weaken browser security or install an untrusted camera workaround.

A production Observer Mode requires one of:

\begin{enumerate}[leftmargin=*]
  \item local HTTPS with a certificate chain trusted by the phone; or
  \item a signed native Pilot application with explicit camera permission,
        foreground indicators, operating-system lifecycle handling, and a
        reviewed privacy declaration.
\end{enumerate}

Until one of these paths exists, repeated automatic camera observation is a
backend-complete but browser-constrained capability; manual owner capture is
the working developer-preview path.

\section{Capability and Interface Boundaries}

\begin{table}[htbp]
\centering
\small
\caption{Relevant signed Fabric capabilities.}
\begin{tabularx}{\textwidth}{@{}L{58mm}L{14mm}Y@{}}
\toprule
Capability & Risk & Scope\\
\midrule
\nolinkurl{voice.live.fact_check} & Low & Build explicit live context and request
evidence-labelled verification.\\
\nolinkurl{companion.webos.perception.submit} & Low & Submit one explicit bounded
phone image for local ephemeral analysis.\\
\nolinkurl{companion.screen.context.correct} & Low & Correct one allowlisted field
of the current structured observation.\\
\nolinkurl{companion.screen.observer.control} & Medium & Start or stop one visible,
short-lived, persona-bound observer session.\\
\bottomrule
\end{tabularx}
\end{table}

Phone operations use a rotating private companion path and a separate CSRF
secret.  Perception requests are limited to 5 MiB.  Ordinary controller
requests are limited separately.  Observer frames carry the additional bearer
token.  The server returns permission failures distinctly from malformed input
and unexpected processing failures.

The phone response includes a restrictive camera policy scoped to the current
origin and disables microphone access on the controller page.  The camera
belongs to the visual-observer workflow; television audio continues to be
handled by the separately governed local voice service.

\section{Data Lifecycle and Privacy Model}

\begin{table}[htbp]
\centering
\small
\caption{Retention by data class.}
\begin{tabularx}{\textwidth}{@{}L{35mm}L{35mm}Y@{}}
\toprule
Data & Retention & Persisted derivative\\
\midrule
Recent microphone audio & Memory-only; cleared on sleep or stop & Bounded
transcript fragments only after explicit contextual request\\
Phone photograph & Temporary local file deleted after analysis & Hash, byte
count, OCR, geometry, labels, and inference\\
Observer frame & Same immediate deletion as manual capture & Structured
perception, fusion result, hash, and session counter\\
Observer bearer token & Never persisted in raw form & SHA-256 token hash while
session remains usable\\
Provider response & Raw response not retained & Bounded normalized metadata
fields and evidence hash\\
Owner correction & Persisted, persona-bound, maximum 240 characters & Field,
value, timestamp, persona ID, and updated evidence hash\\
Pilot Moment & Persisted signed context record & No protected audio or video\\
\bottomrule
\end{tabularx}
\end{table}

The system does not infer face identity.  A shared television request cannot
select a private payment method, account, contact, or calendar automatically.
Products, locations, purchases, bookings, and sharing remain within the
existing private-phone confirmation model.

\section{Implementation Map}

\begin{table}[htbp]
\centering
\small
\begin{tabularx}{\textwidth}{@{}L{43mm}Y@{}}
\toprule
Module & Responsibility\\
\midrule
\texttt{voice.py} & Wake phrase, microphone lifecycle, memory-only audio ring,
bounded transcript context, and live intents.\\
\texttt{live\_context.py} & Construction and canonical hashing of explicit
LiveContext evidence.\\
\texttt{moments.py} & Rights-aware Moment payload, canonical hash, device
signature, and prepared-not-sent state.\\
\texttt{apple\_vision.swift} & Local text recognition, normalized geometry, and
image classifications.\\
\texttt{perception.py} & Input validation, temporary-file lifecycle, structured
perception record, and immediate image deletion.\\
\texttt{screen\_fusion.py} & Deterministic source-aware fusion, supported
detectors, confidence, conflict state, correction, and evidence hash.\\
\texttt{provider\_metadata.py} & Provider URL normalization and optional
bounded official YouTube metadata enrichment.\\
\texttt{observer.py} & Explicit session lease, persona binding, token hashing,
expiry, frame interval enforcement, and stop.\\
\texttt{live\_news.py} & Verdict normalization, per-claim evidence binding,
source classification, correction timeline, privacy declaration, canonical
hash, and bounded history.\\
\texttt{companion.py} & Private phone panels, manual capture, Observer controls,
camera preview, CSRF, request limits, and frame upload.\\
\texttt{runtime.py} & Signed capability enforcement, evidence orchestration,
provider fusion, audit events, and answer routing.\\
\texttt{cli.py} & Dashboard, phone snapshot projection, controller services,
and runtime status.\\
\bottomrule
\end{tabularx}
\end{table}

\section{Verification and Reproducibility}

Automated verification covers provider parsing, official-API normalization,
entitlement non-claims, Observer token redaction, invalid-token rejection,
frame-rate enforcement, persona-bound stop, screen fusion, subtitle and score
detection, product confirmation boundaries, game-HUD detection, correction
allowlists, companion controls, and the existing Fabric regression suite.

The terminal release evidence is:

\begin{itemize}[leftmargin=*]
  \item Fabric 0.15.0: 136 passing checks; package SHA-256
        \nolinkurl{c313bdd26da53f1db452e34377d2aef6bdc6dc2510822a64cb54ac439c3d1ed3}.
  \item Fabric 0.16.0: 144 passing checks; package SHA-256
        \nolinkurl{5304d4bdd1beb724fd96f3ee90a14983e7a67021f51e333b3ac38fbf87aff595}.
  \item Fabric 0.17.0: 149 passing checks; package SHA-256
        \nolinkurl{513eac8a8c0c9e4178837a16103c4266156341140e609af6430dc920054d49a1}.
  \item Fabric 0.18.0: 154 passing checks; package SHA-256
        \nolinkurl{ac594129b59a5a2a97110ace96a2aed03013ee4164db0a79026c0641d7ad0eae}.
  \item Fabric 0.18.1: 156 passing checks; package SHA-256
        \nolinkurl{897b15e348d9367a633581ac4b4a12a6a4063ad095e924757e0f2c737312afd0}.
\end{itemize}

Fabric 0.15 was field-verified against the paired LG television for permitted
state observation and Moment preparation.  Fabric 0.16 verified real Apple
Vision OCR geometry.  The 0.16 and 0.17 implementation and regression runs did
not issue live TV control actions, allowing the household game session to
continue uninterrupted.  At the 0.17 release gate the audit chain was valid,
the microphone was physically closed, and Observer Mode was stopped.

\section{Fabric 0.18: Claim-Decomposed Live News}

Fabric 0.18 turns the earlier evidence-labelled fact-check proof into a
production-shaped workflow.  The structured research contract requires an
overall verdict and confidence, at most four independently checkable claims,
per-claim verdicts and explanations, explicit evidence indexes, date scope,
context notes, and at most five HTTPS evidence sources.  The provider request
uses built-in web search and strict JSON Schema output with response storage
disabled.  Fabric retains only the sanitized structured result.

The mother then compiles \texttt{pilot.live-news.v1}.  Evidence indexes that do
not correspond to retained HTTPS sources are removed.  Missing provider
authority forces \texttt{Unverifiable} with zero confidence.  The evidence
record includes a chronological statement-capture, owner-correction, and
research timeline plus a canonical digest.  Raw audio, source pixels, raw
provider responses, and private identity are not projected or retained in the
fact-check object.

The presentation model is intentionally asymmetric.  Pilot speaks the compact
verdict without navigating away from the active programme and exposes the full
claim/source/timeline record on the private phone.  ``Pilot, show me the
evidence'' is an explicit request to open the full-screen Canvas.  A persistent
overlay above arbitrary third-party webOS applications is not claimed: LG
documents notification overlays as system-application functionality.  A true
same-screen overlay therefore remains native-platform work rather than a
gateway simulation.

\section{Honest Capability Boundary}

The implemented proof establishes that Pilot can:

\begin{itemize}[leftmargin=*]
  \item preserve a bounded recent context without retaining raw room audio;
  \item create a signed current-programme context after explicit request;
  \item research and display evidence-labelled explanations or fact checks;
  \item prepare a rights-aware, signed Moment without copying protected media;
  \item extract OCR text and geometry locally from an explicit phone image;
  \item recognize a bounded set of subtitle, score, product, location, and game
        cues;
  \item fuse independent sources while retaining source and confidence labels;
  \item let a private persona correct the current observation;
  \item normalize supported provider identifiers and optionally enrich YouTube
        metadata; and
  \item create a visible, expiring, rate-limited Observer authority.
\end{itemize}

The proof does \emph{not} yet establish that Pilot can:

\begin{itemize}[leftmargin=*]
  \item see arbitrary protected TV pixels without the owner-provided camera;
  \item run automatic phone-camera observation over insecure LAN HTTP;
  \item identify an actor or household member solely from a face;
  \item guarantee that OCR text is an exact subtitle or factual statement;
  \item establish a household's subscription or playback entitlement from a
        provider URL;
  \item create or distribute a protected media clip where the provider has not
        authorized it;
  \item send a Moment without separate private recipient selection and
        authorization; or
  \item provide production-grade continuous perception without native/TLS
        lifecycle, battery, thermal, privacy, and accessibility testing.
\end{itemize}

\section{Fabric 0.19: AION Native Intelligence Routing}

Fabric 0.19 removes OpenAI from the default research path.  The mother brain
stores one of three maximum-authority policies: \texttt{native},
\texttt{gemini}, or \texttt{boost}.  In every policy the route begins with
deterministic local capability, available AION state and memory, the installed
local model, and bounded public evidence.  The Gemini policy may then use a
persona-provided Gemini key for evidence synthesis; Google Search grounding is
a separate, explicitly enabled capability.  The Boost policy may reach a
premium provider only after the preceding routes fail.

The routing record deliberately contains no secret.  It records a provider
class trace, local UTC usage counts, and one user-facing provenance label:
\texttt{AION Local}, \texttt{AION + live public evidence}, or
\texttt{Pilot Boost}.  Gemini requests are constrained by a configurable local
daily safety allowance.  The provider key remains in the mother-brain vault and
is never projected to the television or edge nodes.

Evidence checking is not delegated entirely to a model.  Fabric filters to
bounded HTTPS sources, rejects model-introduced URLs when an allowed evidence
set exists, validates claim citation indexes, and detects contradictory
verdicts for normalized duplicate claims.  In the internet/no-paid-AI state,
Pilot can still present useful current sources.  In fact-check mode it returns
\texttt{Unverifiable} when no model result passes the boundary instead of
manufacturing a verdict.

\section{Remaining Production Work}

The next production-quality steps are:

\begin{enumerate}[leftmargin=*]
  \item deploy trusted local TLS or a signed native phone application;
  \item add a persistent foreground camera indicator, explicit pause, session
        countdown, battery/thermal adaptation, and operating-system suspension
        recovery;
  \item connect authenticated, persona-specific provider metadata and
        entitlement adapters without exposing account credentials to the TV;
  \item add official provider-native timestamp, Moment, and clip execution
        where terms and APIs permit it;
  \item implement observe--act--observe--verify for supported third-party
        navigation while preventing automatic legal consent or purchases;
  \item improve subtitle, scoreboard, product, and programme identification
        through provider metadata and temporally stable multi-frame fusion;
  \item add spoiler-aware explanation, translation, sports interpretation, and
        native-platform compact overlays where the TV platform grants authority; and
  \item complete privacy, accessibility, child-safety, rights, abuse, and
        regional compliance review before consumer release.
\end{enumerate}

\section{Conclusion}

Fabric 0.15--0.19 establishes the first governed perception loop for Pilot's
television product.  The important architectural result is not a claim that the
restricted TV has become an unrestricted camera or media extractor.  It is
that the mother brain can combine explicit, bounded, independently labelled
evidence into useful room-scale intelligence while preserving ownership,
privacy, rights, and action boundaries.

The resulting pattern is:

\begin{center}
\textbf{recent context $+$ permitted device state $+$ owner-visible pixels
$+$ provider evidence}\\[4pt]
$\downarrow$\\[4pt]
\textbf{confidence-labelled understanding $\rightarrow$ private correction
$\rightarrow$ governed answer or action.}
\end{center}

This is the foundation for live fact checking, contextual explanation,
translation, sports intelligence, product discovery, accessibility, education,
and rights-aware sharing without requiring a replacement television.

\begin{thebibliography}{9}

\bibitem{youtubevideos}
Google for Developers, ``YouTube Data API: Videos: list.''\\
\url{https://developers.google.com/youtube/v3/docs/videos/list}

\bibitem{netflixmoments}
Netflix Help Center, ``How to save and share clips from Netflix.''\\
\url{https://help.netflix.com/en/node/210664027435620}

\bibitem{twitchembed}
Twitch Developers, ``Embedding Video and Clips'' and ``Everything.''\\
\url{https://dev.twitch.tv/docs/embed/video-and-clips/}\\
\url{https://dev.twitch.tv/docs/embed/}

\bibitem{mdnmedia}
MDN Web Docs, ``MediaDevices: getUserMedia() method.''\\
\url{https://developer.mozilla.org/en-US/docs/Web/API/MediaDevices/getUserMedia}

\bibitem{webosscreenshot}
LG webOS TV Developer Forum, ``Is there an API like in Android TV get
screenshot of the screen?''\\
\url{https://forum.webostv.developer.lge.com/t/is-there-an-api-like-in-android-tv-get-screenshot-of-the-screen/1419}

\bibitem{openairesponses}
OpenAI, ``Responses API: Create a model response.''\\
\url{https://developers.openai.com/api/reference/resources/responses/methods/create}

\bibitem{webosnotifications}
LG webOS TV Developer, ``Notifications.''\\
\url{https://webostv.developer.lge.com/develop/guides/notifications}

\end{thebibliography}

\end{document}
